% !TEX program = lualatex
% Generated by new_lecture_note.sh
\documentclass[11pt]{article}

\usepackage{fontspec}
\usepackage{microtype}
\usepackage{mathtools,amssymb,amsfonts,amsthm,bm}
\usepackage{enumitem}
\usepackage{booktabs,array,xcolor}
\usepackage{geometry,fancyhdr,setspace}
\usepackage[hidelinks]{hyperref}
\usepackage[nameinlink,noabbrev]{cleveref}

\geometry{margin=1in,headheight=15pt}
\setstretch{1.12}
\setlength{\parindent}{0pt}
\setlength{\parskip}{0.55em plus 0.12em minus 0.08em}
\setlength{\abovedisplayskip}{0.8em plus 0.2em minus 0.1em}
\setlength{\belowdisplayskip}{0.8em plus 0.2em minus 0.1em}
\allowdisplaybreaks[3]
\setlist{leftmargin=*,topsep=0.35em,itemsep=0.15em,parsep=0pt}

% ---------- Metadata ----------
\newcommand{\studentname}{Keshav Ramamurthy}
\newcommand{\coursename}{Stat 150}
\newcommand{\courselongname}{Stochastic Processes}
\newcommand{\lecturenumber}{02}
\newcommand{\lecturetitle}{Lecture 02}
\newcommand{\lecturedate}{\today}

% ---------- Reusable notation ----------
\newcommand{\N}{\mathbb{N}}
\newcommand{\Z}{\mathbb{Z}}
\newcommand{\Q}{\mathbb{Q}}
\newcommand{\R}{\mathbb{R}}
\newcommand{\C}{\mathbb{C}}
\newcommand{\F}{\mathbb{F}}
\newcommand{\K}{\mathbb{K}}
\newcommand{\E}{\mathbb{E}}
\newcommand{\Prob}{\mathbb{P}}
\newcommand{\1}{\mathbf{1}}
\newcommand{\eps}{\varepsilon}
\newcommand{\dd}{\,\mathrm{d}}
\newcommand{\abs}[1]{\left\lvert#1\right\rvert}
\newcommand{\norm}[1]{\left\lVert#1\right\rVert}
\newcommand{\ip}[2]{\left\langle#1,#2\right\rangle}
\newcommand{\set}[1]{\left\{#1\right\}}
\newcommand{\given}{\,\middle\vert\,}
\newcommand{\indep}{\perp\!\!\!\perp}
\newcommand{\lto}{\longrightarrow}
\newcommand{\deriv}[2]{\frac{\mathrm{d}#1}{\mathrm{d}#2}}
\newcommand{\pderiv}[2]{\frac{\partial#1}{\partial#2}}
\DeclareMathOperator{\Aut}{Aut}
\DeclareMathOperator{\Hom}{Hom}
\DeclareMathOperator{\Ker}{ker}
\DeclareMathOperator{\im}{im}
\DeclareMathOperator{\rank}{rank}
\DeclareMathOperator{\nullity}{nullity}
\DeclareMathOperator{\Span}{span}
\DeclareMathOperator{\Var}{Var}
\DeclareMathOperator{\Cov}{Cov}

% ---------- Note environments ----------
\newtheoremstyle{notestyle}{0.7em}{0.7em}{\itshape}{}{}{\bfseries}{.5em}{}
\theoremstyle{notestyle}
\newtheorem{theorem}{Theorem}[section]
\newtheorem{lemma}[theorem]{Lemma}
\newtheorem{proposition}[theorem]{Proposition}
\newtheorem{corollary}[theorem]{Corollary}
\theoremstyle{definition}
\newtheorem{definition}[theorem]{Definition}
\newtheorem{example}[theorem]{Example}
\newtheorem{remark}[theorem]{Remark}
\newenvironment{claim}{\par\noindent\textbf{Claim.}\ }{\par}
\newenvironment{idea}{\par\noindent\textit{Idea.}\ }{\par}
\newenvironment{proofsketch}{\par\noindent\textbf{Proof sketch.}\ }{\par}

% ---------- Header ----------
\pagestyle{fancy}
\fancyhf{}
\fancyhead[L]{\coursename}
\fancyhead[C]{\lecturetitle}
\fancyhead[R]{\studentname}
\fancyfoot[C]{\thepage}

\title{\vspace{-1.5em}\textbf{\coursename\ -- \lecturetitle}}
\author{\studentname}
\date{\lecturedate}

\begin{document}
\maketitle

\section{Markov Chain 2.0}
\subsection{Classification of States:}
Understand the behavior of Markov chains. Sometimes they are recurrent, sometimes
transient, etc.
\par Notation: $i \rightarrow j$ denotes ``$j$ is reachable or accessible from $i$'',
that is
$$\exists n\in \mathbb{Z}^{+} \text{ such that } P_{ij}^{n} = \text{
probability of jump to j at n steps starting from i} > 0$$
\par Write $i \leftrightarrow j$ if $i \rightarrow j$ \emph{and} $j \rightarrow i$
(they \textbf{communicate}). Communication is an equivalence relation, so it splits
the state space into \textbf{communication classes}; the chain is
\textbf{irreducible} if there is only one class, the whole state space.
\begin{definition}[Chapman--Kolmogorov Equation -- Discrete Case]
\label{def:label}
To travel from $i$ to $k$ in $m + n$ steps, the chain passes through some
intermediate state $r$ at time $m$; summing over all intermediates,
$$P_{ik}^{m + n} = \sum_{r} P_{ir}^{m}P_{rk}^{n}$$
In matrix form this is exactly $P^{m+n} = P^{m} P^{n}$: multi-step transition
matrices compose by ordinary matrix multiplication.
\end{definition}
Notation:
\begin{enumerate}
    \item $\mathbb{P}_x(A) = \mathbb P(a \mid x_0 = x)$ for an event $A$ in a markov
        chain $(x_n)_{n \in \mathbb{N}}$
    \item Denote $T_y = \min\{n \ge 1 : X_n = y\}$  time of the first return to y. (y at
        time 0, return back).
    \item $$\rho_{yy} = \mathbb{P}_{y}(T_y < \infty) =
    \sum_{n=1}^{\infty}\mathbb{P}_y(T_y = n)$$  is the probability that, starting at
    y, we return in a finite number of steps. (Durrett writes this $\rho_{yy}$,
    ``rho'', rather than $P_{yy}$, to avoid confusion with the transition
    probabilities $P(i,j)$.)
\end{enumerate}
\begin{definition}[Recurrent and Transient States]
\label{def:recurrent-transient}
A state $y$ is \textbf{recurrent} if $\rho_{yy} = 1$ and \textbf{transient} if
$\rho_{yy} < 1$.
\end{definition}
\par Why this dichotomy is the whole story: by the strong Markov property, each
return to $y$ ``restarts'' the chain, so successive returns behave like
independent trials with success probability $\rho_{yy}$. The probability of a
$k$th return is $(\rho_{yy})^{k}$, the probability of exactly $k$ returns is
$(\rho_{yy})^{k}(1 - \rho_{yy})$, and the expected number of returns is
$\rho_{yy}/(1 - \rho_{yy})$. The chain returns to $y$ infinitely often with
probability 1 \emph{exactly when} $\rho_{yy} = 1$ -- there is no middle ground.
(This is Durrett \S1.3, and it is where the ``kth return is $(\rho_{yy})^{k}$''
claim from lecture comes from.)
\begin{definition}[Strong Markov Property]
\label{def:label}
The memoryless property survives stopping at \emph{random} times: stopped at a
stopping time and restarted, the chain is again a Markov chain with the same
transitions, independent of the past. This is what powers results like ``the
probability of a $k$th return is $(\rho_{yy})^{k}$.''
\end{definition}
\begin{definition}[Stopping Time]
\label{def:label}
A random time $T$ (often written $\tau$) is a stopping time if for every $n$ the
event $\{\tau = n\}$ (``we stop at time $n$'') is determined by
$X_0, X_1, \cdots, X_n$ only -- no peeking at the future. Examples: the first
return time $T_y$; the first time the chain enters a set $A$ of states.
\end{definition}
\newpage
\textbf{1.2 Durrett: Strong Markov Chain Properties}
Let $(X_t)_{t \in \mathbb{N}}$ be a Markov chain, $\tau$ a stopping time, and $E$ an
event determined by the history up to $\tau$. Then for all $i, j \in S$ and $k \ge 1$,
$$\mathbb{P}(X_{\tau + k} = j \mid X_{\tau} = i, E) = \mathbb{P}(X_{\tau + k} = j \mid X_{\tau} = i),$$
and the shifted chain $(Y_n) \coloneqq X_{\tau + n}$, $n \ge 0$, is itself a Markov
chain with the same transition matrix.
\par \textbf{Sketch of Proof: } condition on $\{\tau = n\}$ (allowed, because
$\{\tau = n\}$ is determined by $X_0, \dots, X_n$) and apply the ordinary Markov
property. We eventually show
$$\mathbb{P}(X_{\tau + 1} = z \mid X_{\tau} = y, \tau = n) = P_{yz} = P(y,z),$$
then sum over $n$ to drop the conditioning on the particular stopping time.

\section{Questions and follow-up}
\begin{itemize}
  \item Connect the lectures: for the symmetric walk on $\mathbb{Z}$ (recurrent),
    what is $\rho_{yy}$, and why can the expected return time still be infinite?
  \item Redo Durrett's \S1.3 computation showing transience when $p \neq
    \frac{1}{2}$.
  \item TODO:
\end{itemize}

\end{document}
